\documentclass[12pt]{article}

\usepackage[margin=1in]{geometry}
\usepackage{amsmath,amssymb,amsthm}
\usepackage{booktabs}
\usepackage{hyperref}
\usepackage{enumitem}

% Theorem environments
\newtheorem{theorem}{Theorem}[section]
\newtheorem{lemma}[theorem]{Lemma}
\newtheorem{proposition}[theorem]{Proposition}
\newtheorem{corollary}[theorem]{Corollary}
\newtheorem{conjecture}[theorem]{Conjecture}
\theoremstyle{definition}
\newtheorem{definition}[theorem]{Definition}
\newtheorem{example}[theorem]{Example}
\theoremstyle{remark}
\newtheorem{remark}[theorem]{Remark}

% Notation
\newcommand{\Z}{\mathbb{Z}}
\newcommand{\Zk}{\Z/k\Z}
\newcommand{\GZk}{G(\Z_k)}

\title{On the Chromatic Number of Central Digraphs over Cyclic Groups}

\author{
% Authors
}

\date{}

\begin{document}

\maketitle

\begin{abstract}
We study the chromatic number of the underlying undirected graph of the
natural central digraph $G(k,1,1)$ on $\Zk \times \Zk$ for $k = 3$ to~$13$.
We prove that for all $k \ge 3$ the graph has diameter exactly~$2$,
degree sequence $(2k{-}2, 2k{-}1)$, edge count $k(k{-}1)(2k{+}1)/2$,
and clique number $\omega = 3$.  Computationally we establish the chromatic
number $\chi(\GZk)$ for each $k$ in this range, revealing a
\emph{staircase} growth pattern: $\chi = 4$ for $k = 3$; $\chi = 5$ for
$k = 4,5,6$; $\chi = 6$ for $k = 7,8,9,10$; and $\chi = 7$ for
$k = 11,12,13$.  This refutes the earlier conjecture that $\chi = 5$
for all $k \ge 4$.  We also record closed-form expressions for the triangle
count and independence number, and give partial chromatic polynomials for
$k = 3$ and~$4$.
\end{abstract}


%======================================================================
\section{Introduction}
\label{sec:intro}
%======================================================================

Central digraphs, introduced in the study of semigroup actions on finite
sets, give rise to highly structured directed graphs whose colouring
properties are not well understood.  For a positive integer $k \ge 3$, the
\emph{natural central digraph} $G(k,1,1)$ on $\Zk \times \Zk$ has vertex
set $\Zk \times \Zk$ and arcs
\[
  (p,\, i) \;\longrightarrow\; \bigl((p+i) \bmod k,\; j\bigr)
  \qquad \text{for every } j \in \Zk.
\]
We study the \emph{underlying undirected graph}, obtained by ignoring arc
directions and removing self-loops.  Throughout, ``colouring'' means a
proper vertex colouring of this undirected graph, and $\chi(\GZk)$ denotes
its chromatic number.

Early computations for small~$k$ led to the conjecture that $\chi(\GZk) = 5$
for all $k \ge 4$ (see Rawat \& Sehgal~\cite{rawat2026}).  We show this is
false: $\chi(\GZk)$ grows in a staircase pattern with jumps at $k = 4, 7,$
and~$11$.

\medskip
\noindent\textbf{Main results.}
\begin{enumerate}[label=(\roman*)]
  \item \emph{Structural invariants (all $k \ge 3$, proved):}
    $\operatorname{diam}(\GZk) = 2$, $\;\omega(\GZk) = 3$,
    $\;|E| = k(k-1)(2k+1)/2$, and every vertex has degree $2k-2$ or $2k-1$.
  \item \emph{Chromatic number (proved for $k \le 10$, strong evidence for
    $k = 11$--$13$):}
    \[
      \chi(\GZk) =
      \begin{cases}
        4 & k = 3,\\
        5 & k = 4, 5, 6,\\
        6 & k = 7, 8, 9, 10,\\
        7 & k = 11, 12, 13.
      \end{cases}
    \]
  \item \emph{Independence number (verified computationally for $3 \le k \le 13$):}
    $\alpha(\GZk) = \lfloor k^2/4 \rfloor$.
\end{enumerate}


%======================================================================
\section{Definitions and Notation}
\label{sec:defs}
%======================================================================

\begin{definition}
\label{def:graph}
For an integer $k \ge 3$, the \emph{central digraph} $\GZk$ has
\begin{itemize}
  \item \textbf{Vertex set:} $V = \Zk \times \Zk$, so $|V| = k^2$.
  \item \textbf{Arc set:} $(p, i) \to \bigl((p+i) \bmod k,\; j\bigr)$
    for all $j \in \Zk$.
\end{itemize}
The \emph{underlying undirected graph}, also denoted $\GZk$, has the same
vertex set; distinct vertices $u$ and $v$ are adjacent (written $u \sim v$)
whenever there is an arc $u \to v$ or $v \to u$ (or both) in the digraph.
Self-loops $(p,0) \to (p,0)$ are removed.
\end{definition}

\begin{remark}
\label{rem:selfloops}
The digraph has a self-loop at $(p, i)$ precisely when $p + i \equiv p$
and $i$ is a valid second coordinate, i.e.\ $i = 0$.  So the self-loop
vertices are exactly $\{(p, 0) : p \in \Zk\}$, a set of size~$k$.
\end{remark}


%======================================================================
\section{Structural Invariants}
\label{sec:structure}
%======================================================================

\subsection{Adjacency characterisation}

\begin{lemma}[Adjacency criterion]
\label{lem:adj}
Two distinct vertices $(p_1, i_1)$ and $(p_2, i_2)$ are adjacent in $\GZk$
if and only if
\[
  (p_2 - p_1) \bmod k \;\in\; \{i_1,\; (-i_2) \bmod k\}.
\]
\end{lemma}

\begin{proof}
Write $\delta = (p_2 - p_1) \bmod k$.  There is an arc $(p_1, i_1) \to (p_2, i_2)$
precisely when $p_2 \equiv p_1 + i_1 \pmod{k}$, i.e.\ $\delta = i_1$.
There is an arc $(p_2, i_2) \to (p_1, i_1)$ precisely when
$p_1 \equiv p_2 + i_2 \pmod{k}$, i.e.\ $\delta \equiv -i_2 \pmod{k}$.
The undirected edge exists if and only if at least one arc is present.
\end{proof}

\subsection{Degree sequence}

\begin{proposition}[Degrees]
\label{prop:deg}
For every $k \ge 3$:
\begin{enumerate}[label=(\alph*)]
  \item Each vertex $(p, 0)$ has degree $2k - 2$.
  \item Each vertex $(p, i)$ with $i \ne 0$ has degree $2k - 1$.
\end{enumerate}
\end{proposition}

\begin{proof}
The out-neighbourhood of $(p, i)$ in the digraph is
$N^+(p,i) = \{((p+i) \bmod k,\; j) : j \in \Zk\}$,
the full ``column'' with first coordinate $(p+i) \bmod k$.
The in-neighbourhood is
$N^-(p,i) = \{((p-j) \bmod k,\; j) : j \in \Zk\}$.

\smallskip\noindent
\textbf{Overlap.} A vertex $w$ lies in $N^+(p,i) \cap N^-(p,i)$ exactly
when $w = ((p+i) \bmod k,\; j)$ and $w = ((p-j) \bmod k,\; j)$, giving
$p + i \equiv p - j \pmod{k}$, i.e.\ $j \equiv -i \pmod{k}$.  The unique
overlap vertex is $w_0 = ((p+i) \bmod k,\; (-i) \bmod k)$.

\smallskip\noindent
\textbf{Self-loops.}  Vertex $(p,i)$ belongs to $N^+(p,i)$ iff $i = 0$
(its column is then $p$, the same as its own first coordinate).
Vertex $(p,i)$ belongs to $N^-(p,i)$ iff $i = 0$ (set $j = i$ in the
in-neighbourhood formula: $((p - i) \bmod k, i) = (p, i)$ requires $i = 0$).

\smallskip\noindent
\textbf{Case $i = 0$.}  Both $N^+$ and $N^-$ contain $(p, 0)$ itself,
so the undirected neighbourhood has size
$(|N^+| - 1) + (|N^-| - 1) - (|\text{overlap}| - 1) = (k-1) + (k-1) - 0 = 2k-2$,
where we subtract the self-loop from each set and note the overlap vertex
$w_0 = (p, 0)$ is the self-loop itself, already removed.

\smallskip\noindent
\textbf{Case $i \ne 0$.}  Neither $N^+$ nor $N^-$ contains $(p,i)$, since
$p + i \not\equiv p$ and the in-neighbourhood argument above also fails.
The overlap vertex $w_0 \ne (p, i)$ is a genuine common neighbour.
The undirected neighbourhood has size
$|N^+| + |N^-| - |\text{overlap}| = k + k - 1 = 2k - 1$.
\end{proof}


\subsection{Edge count}

\begin{theorem}[Edge count]
\label{thm:edges}
For every $k \ge 3$,
\[
  |E(\GZk)| \;=\; \frac{k(k-1)(2k+1)}{2}.
\]
\end{theorem}

\begin{proof}
By the handshaking lemma, $2|E| = \sum_{v \in V} \deg(v)$.
There are $k$ vertices of the form $(p, 0)$, each of degree $2k - 2$,
and $k(k-1)$ vertices $(p, i)$ with $i \ne 0$, each of degree $2k - 1$.
Therefore
\[
  2|E| = k(2k - 2) + k(k - 1)(2k - 1)
       = k(k-1)\bigl[2 + (2k - 1)\bigr]
       = k(k-1)(2k+1),
\]
and the result follows.
\end{proof}


\subsection{Connectivity and diameter}

\begin{theorem}[Diameter]
\label{thm:diam}
For every $k \ge 3$, the graph $\GZk$ is connected with
$\operatorname{diam}(\GZk) = 2$.
\end{theorem}

\begin{proof}
We show that every pair of distinct vertices is at distance at most~$2$,
and exhibit a pair at distance exactly~$2$.

\smallskip\noindent\textbf{Upper bound.}
Let $u = (p_1, i_1)$ and $v = (p_2, i_2)$ be distinct vertices.
Define $c_1 = (p_1 + i_1) \bmod k$ and set
\[
  w \;=\; \bigl(c_1,\; (p_2 - c_1) \bmod k\bigr).
\]
We claim $w$ is a common neighbour of $u$ and $v$ (and is distinct from both
whenever $u \not\sim v$).

\emph{$u \sim w$:}  The first coordinate of $w$ is $c_1 = (p_1 + i_1) \bmod k$,
so $w$ lies in the out-neighbourhood of~$u$.  Since $u$ is adjacent to every
vertex in this column, $u \sim w$.

\emph{$v \sim w$:}  The vertex $w$ sends an arc to column
$(c_1 + (p_2 - c_1)) \bmod k = p_2$.  Since $v = (p_2, i_2)$ lies in column~$p_2$,
the arc $w \to v$ exists, so $w \sim v$.

\emph{Distinctness:}  If $w = u$, then $c_1 = p_1$ (giving $i_1 = 0$) and
$p_2 - c_1 \equiv i_1 = 0$, so $p_2 = p_1$.  Then $u = (p_1, 0)$ and
$v = (p_1, i_2)$ are in the same column with $u$ having second coordinate~$0$,
whence $u \sim v$ directly.  Similarly, $w = v$ forces $c_1 = p_2$ and $i_2 = 0$,
so $u \to v$ and again $u \sim v$.  In either degenerate case $d(u,v) = 1 \le 2$.

\smallskip\noindent\textbf{Lower bound.}
Take $u = (1, 2)$ and $v = (2, 1)$.  We have $\delta = 1$ and
$\{i_1,\, (-i_2) \bmod k\} = \{2,\, k-1\}$.  For $k \ge 3$,
$1 \notin \{2, k-1\}$, so $u \not\sim v$ and $d(u,v) \ge 2$.
\end{proof}


\subsection{Clique number}

\begin{theorem}[Clique number]
\label{thm:clique}
For every $k \ge 3$, $\;\omega(\GZk) = 3$.
\end{theorem}

\begin{proof}
\emph{Lower bound ($\omega \ge 3$).}
The vertices $(0,1)$, $(1,0)$, and $(1,1)$ form a triangle:
\begin{align*}
  (0,1) \sim (1,0) &:\; \delta = 1 \in \{1, 0\}, \\
  (0,1) \sim (1,1) &:\; \delta = 1 \in \{1, k{-}1\}, \\
  (1,0) \sim (1,1) &:\; \delta = 0 \in \{0, k{-}1\}.
\end{align*}

\smallskip
\emph{Upper bound ($\omega \le 3$).}
Suppose for contradiction that $\GZk$ contains a clique
$K = \{v_1, v_2, v_3, v_4\}$ of size~$4$, where $v_a = (p_a, i_a)$.
We use the following structural observation.

\begin{quote}
\textbf{Column constraint.}  All digraph arcs from vertex $(p,i)$
terminate at vertices whose first coordinate is $(p + i) \bmod k$.
Therefore, among the other three clique members, those reached by a
\emph{forward arc} from~$v_a$ (i.e.\ an arc $v_a \to v_b$, equivalently
$p_b \equiv p_a + i_a$) all share the same first coordinate.
\end{quote}

We say $v_a$ \emph{arcs to} $v_b$ if there is a digraph arc $v_a \to v_b$.
Each undirected edge $\{v_a, v_b\}$ requires at least one arc between
$v_a$ and $v_b$ (Lemma~\ref{lem:adj}), which counts as a forward arc
from one endpoint.

\smallskip\noindent
\textbf{Step 1: No vertex arcs to all three others.}
Suppose $v_1$ arcs to $v_2, v_3, v_4$.  Then $p_2 = p_3 = p_4 =: c$.
The edge $v_2 \sim v_3$ requires $\delta_{23} = 0 \in \{i_2, -i_3\}$,
so $i_2 = 0$ or $i_3 = 0$.  Similarly $v_2 \sim v_4$ gives $i_2 = 0$ or
$i_4 = 0$, and $v_3 \sim v_4$ gives $i_3 = 0$ or $i_4 = 0$.
If $i_2 = 0$, then $v_3 \sim v_4$ forces $i_3 = 0$ or $i_4 = 0$, but
$i_3 = 0$ gives $v_3 = (c, 0) = v_2$, and $i_4 = 0$ gives
$v_4 = (c, 0) = v_2$ --- contradictions.
The cases $i_3 = 0$ and $i_4 = 0$ are symmetric.

\smallskip\noindent
\textbf{Step 2: No vertex arcs to exactly two others.}
Suppose $v_1$ arcs to $v_2$ and $v_3$ (but not~$v_4$).
Then $p_2 = p_3 =: c = (p_1 + i_1) \bmod k$.  The edge $v_2 \sim v_3$
gives $i_2 = 0$ or $i_3 = 0$ (as above); assume without loss of generality that
$i_2 = 0$, so $v_2 = (c, 0)$.

Since $v_1$ does not arc to $v_4$, the edge $v_1 \sim v_4$ is realised by
the arc $v_4 \to v_1$, giving $p_1 = (p_4 + i_4) \bmod k$.

The edge $v_2 \sim v_4$ requires $(p_4 - c) \bmod k \in \{0, -i_4\}$.

\begin{itemize}
\item If $p_4 \equiv c$: the edge $v_3 \sim v_4$ requires
  $0 \in \{i_3, -i_4\}$.  If $i_3 = 0$ then $v_3 = (c, 0) = v_2$
  (contradiction); if $i_4 = 0$ then $p_1 = p_4 = c$ and
  $c = p_1 + i_1$ forces $i_1 = 0$, whence $v_1 = (c, 0) = v_2$
  (contradiction).

\item If $p_4 \equiv c - i_4$: then $p_1 = p_4 + i_4 = c$.
  But $c = p_1 + i_1$ gives $i_1 = 0$, so $v_1 = (c, 0) = v_2$
  (contradiction).
\end{itemize}

\smallskip\noindent
\textbf{Step 3: Counting argument.}
From Steps~1 and~2, each clique member arcs to at most one other member.
Each such forward arc services exactly one edge.  So the four vertices
produce at most~$4$ forward arcs.  But the clique has $\binom{4}{2} = 6$
edges, each requiring at least one forward arc (in one direction).
Since $4 < 6$, this is a contradiction.
\end{proof}


\subsection{Triangle count}

\begin{conjecture}[Triangle count]
\label{conj:triangles}
For every $k \ge 3$, the number of triangles in $\GZk$ is
\[
  T(k) \;=\; \frac{k(k-1)(4k-5)}{3}.
\]
\end{conjecture}

This formula has been verified computationally for all $k$ from $3$ to~$13$
(Table~\ref{tab:triangles}).

\begin{table}[h]
\centering
\begin{tabular}{@{}rrrrc@{}}
\toprule
$k$ & $T(k)$ computed & $k(k{-}1)(4k{-}5)/3$ & Match \\
\midrule
3  &    14 &    14 & \checkmark \\
4  &    44 &    44 & \checkmark \\
5  &   100 &   100 & \checkmark \\
6  &   190 &   190 & \checkmark \\
7  &   322 &   322 & \checkmark \\
8  &   504 &   504 & \checkmark \\
9  &   744 &   744 & \checkmark \\
10 &  1050 &  1050 & \checkmark \\
11 &  1430 &  1430 & \checkmark \\
12 &  1892 &  1892 & \checkmark \\
13 &  2444 &  2444 & \checkmark \\
\bottomrule
\end{tabular}
\caption{Triangle counts, verified against the conjectured formula.}
\label{tab:triangles}
\end{table}


%======================================================================
\section{Independence Number}
\label{sec:indep}
%======================================================================

\begin{conjecture}[Independence number]
\label{conj:alpha}
For every $k \ge 3$,
\[
  \alpha(\GZk) \;=\; \left\lfloor \frac{k^2}{4} \right\rfloor
  \;=\;
  \begin{cases}
    k^2/4       & \text{if } k \text{ is even},\\
    (k^2 - 1)/4 & \text{if } k \text{ is odd}.
  \end{cases}
\]
\end{conjecture}

This has been verified computationally for $k = 3$ to~$13$
(Table~\ref{tab:main}).  The fractional relaxation bound
$\chi \ge n/\alpha \ge \lceil 4 \rceil = 4$ (for even~$k$) and
$\lceil k^2 / \lfloor k^2/4 \rfloor \rceil \le 5$ (for all~$k$ tested)
is far below the true chromatic number from $k = 7$ onward, so the
independence number alone does not explain the growth of~$\chi$.


%======================================================================
\section{Chromatic Number}
\label{sec:chromatic}
%======================================================================

\subsection{Summary of results}

\begin{table}[h]
\centering
\begin{tabular}{@{}rrrccccr@{}}
\toprule
$k$ & $|V|$ & $|E|$ & Degrees & $\omega$ & $\alpha$ & $\lceil n/\alpha \rceil$ & $\chi$ \\
\midrule
 3 &   9 &    21 & 4--5   & 3 &  3 & 3 & \textbf{4} \\
 4 &  16 &    54 & 6--7   & 3 &  4 & 4 & \textbf{5} \\
 5 &  25 &   110 & 8--9   & 3 &  6 & 5 & \textbf{5} \\
 6 &  36 &   195 & 10--11 & 3 &  9 & 4 & \textbf{5} \\
 7 &  49 &   315 & 12--13 & 3 & 12 & 5 & \textbf{6} \\
 8 &  64 &   476 & 14--15 & 3 & 16 & 4 & \textbf{6} \\
 9 &  81 &   684 & 16--17 & 3 & 20 & 5 & \textbf{6} \\
10 & 100 &   945 & 18--19 & 3 & 25 & 4 & \textbf{6} \\
11 & 121 & 1265 & 20--21 & 3 & 30 & 5 & \textbf{7} \\
12 & 144 & 1650 & 22--23 & 3 & 36 & 4 & \textbf{7} \\
13 & 169 & 2106 & 24--25 & 3 & 42 & 5 & \textbf{7} \\
\bottomrule
\end{tabular}
\caption{Invariants of $\GZk$ for $k = 3$ to~$13$.  Edge counts and degrees
follow from Theorems~\ref{thm:edges} and Proposition~\ref{prop:deg};
$\omega = 3$ is Theorem~\ref{thm:clique}; $\alpha$ and $\chi$ are
computational (see \S\ref{sec:methods}).}
\label{tab:main}
\end{table}


\subsection{Computational methods}
\label{sec:methods}

Different methods were required as $k$ (and hence $|V| = k^2$) increased.

\subsubsection*{$k = 3$ to $6$: Backtracking}

For $k \le 6$ ($|V| \le 36$), an exhaustive backtracking search over all
proper $K$-colourings determines $\chi$ exactly.  Forward checking prunes
the search at each step: when a vertex is coloured, the available colours
of its uncoloured neighbours are reduced, and the branch is abandoned if
any neighbour's available set becomes empty.

\subsubsection*{$k = 7$ and $8$: DSATUR with symmetry breaking}

At $49$ and $64$ vertices, plain backtracking is too slow.
Two refinements make the search tractable:
\begin{enumerate}
  \item \textbf{DSATUR branching.}  The next vertex to colour is the
    uncoloured vertex whose neighbours already use the most distinct
    colours (the ``saturation degree''); ties go to the vertex of highest
    degree.  This heuristic finds contradictions early.
  \item \textbf{Colour-symmetry breaking.}  Colours are interchangeable,
    so a vertex may use any colour already in play or only the \emph{next}
    unused colour.  This removes up to $K!$ equivalent branches.
\end{enumerate}
Both techniques keep the search exhaustive.  A ``not $K$-colourable''
result is therefore a proof, not a heuristic.

\subsubsection*{$k = 9$ and $10$: SAT encoding}

At $81$ and $100$ vertices the pure-Python DSATUR search becomes too slow.
We encode ``$G$ is $K$-colourable'' as a propositional satisfiability
instance.  For each vertex~$v$ and colour $c \in \{0, \ldots, K{-}1\}$,
introduce a Boolean variable $x_{v,c}$.  The clauses are:
\begin{itemize}
  \item \emph{At least one colour:} $\bigvee_c x_{v,c}$ for every vertex~$v$.
  \item \emph{At most one colour:} $\neg x_{v,c} \lor \neg x_{v,c'}$ for
    every vertex $v$ and $c \ne c'$.
  \item \emph{Edge conflicts:} $\neg x_{u,c} \lor \neg x_{v,c}$ for every
    edge $\{u,v\}$ and every colour~$c$.
  \item \emph{Symmetry breaking:} fix a triangle to colours $0, 1, 2$.
    This is valid because the graph contains a triangle
    (Theorem~\ref{thm:clique}) and colours are interchangeable.
\end{itemize}
The CaDiCaL solver (accessed via the \texttt{python-sat} library) resolves
these instances in seconds.  An UNSAT answer proves no $K$-colouring exists;
a SAT answer returns a colouring, which is verified edge-by-edge by an
independent check.

\subsubsection*{$k = 11$, $12$, and $13$: TabuCol + SAT cross-check}

At $121$+ vertices, neither exhaustive DSATUR nor a direct SAT proof of
non-$6$-colourability finishes within a practical time or conflict budget.
We use a two-pronged strategy.

\medskip\noindent
\textbf{Upper bound ($\chi \le 7$).}
TabuCol, a standard tabu-search metaheuristic for graph colouring
\cite{hertz1987}, finds a proper $7$-colouring in well under $1000$
iterations for each of $k = 11, 12, 13$.  Each $7$-colouring is
independently verified by encoding it as a SAT instance and checking that
CaDiCaL confirms satisfiability in under $0.02$\,s.

\medskip\noindent
\textbf{Lower bound ($\chi \ge 7$).}
Across $20$ independent random seeds and $10{,}000{,}000$ total TabuCol
iterations, no $6$-colouring is found for any of $k = 11, 12, 13$ --- the
search consistently stalls at a small handful of unresolved edge conflicts.
The matching SAT instance for $6$-colourability is also constructed, but
CaDiCaL exhausts tens of millions of conflicts without resolving it either
way.

\medskip
The upper bound $\chi \le 7$ is therefore rigorous (SAT-verified colourings).
The lower bound $\chi \ge 7$ rests on strong heuristic evidence rather than
a formal UNSAT proof.  Turning this into a proof is an open problem
(\S\ref{sec:open}).

\subsubsection*{TabuCol algorithm}

For completeness, we describe the TabuCol procedure used for $k \ge 11$.
\begin{enumerate}
  \item Start from a random (improper) $K$-colouring.  Maintain, for every
    vertex~$v$ and colour~$c$, the count $\gamma(v,c)$ of neighbours of~$v$
    currently coloured~$c$.  A vertex $v$ is \emph{in conflict} when
    $\gamma(v, \operatorname{colour}(v)) > 0$.
  \item At each step, select a conflicted vertex~$v$ and recolour it with
    the colour~$c$ that maximally reduces the total conflict count, breaking
    ties randomly.  A short-term tabu list of length $\ell$ (typically
    $\ell = 7$ to $10$) forbids immediately reversing a move.  An
    \emph{aspiration criterion} overrides the tabu if a move would yield
    fewer total conflicts than the best solution seen so far.
  \item Terminate when the conflict count reaches~$0$ (a proper colouring
    found) or an iteration budget is exhausted.
\end{enumerate}
The $\gamma$ matrix allows each move to be evaluated and updated in $O(\Delta)$
time, where $\Delta$ is the maximum degree.


\subsection{Rigour classification}

Table~\ref{tab:rigour} summarises the proof status of each~$\chi$ value.

\begin{table}[h]
\centering
\begin{tabular}{@{}rcll@{}}
\toprule
$k$ & $\chi$ & Upper bound & Lower bound \\
\midrule
3       & 4 & Exhaustive backtracking & Exhaustive backtracking \\
4--6    & 5 & Exhaustive backtracking & Exhaustive backtracking \\
7--8    & 6 & DSATUR (exhaustive) & DSATUR (exhaustive) \\
9--10   & 6 & SAT (CaDiCaL) & SAT UNSAT (CaDiCaL) \\
11--13  & 7 & SAT-verified colouring & TabuCol (heuristic) \\
\bottomrule
\end{tabular}
\caption{Proof methods for each chromatic number.  All upper bounds are
rigorous.  Lower bounds for $k \le 10$ are proofs; for $k = 11$--$13$
they are strong heuristic evidence.}
\label{tab:rigour}
\end{table}


%======================================================================
\section{The Staircase Pattern}
\label{sec:staircase}
%======================================================================

The chromatic numbers exhibit a striking \emph{staircase} pattern.
Each new value of~$\chi$ first appears at a specific~$k$ and then persists
for a ``plateau'' of consecutive values before the next jump:

\begin{center}
\begin{tabular}{@{}ccccc@{}}
\toprule
$\chi$ & First appears at $k$ & Plateau & Length & Gap from previous \\
\midrule
4 & 3  & $\{3\}$        & 1  & --- \\
5 & 4  & $\{4,5,6\}$    & 3  & $4 - 3 = 1$ (not meaningful) \\
6 & 7  & $\{7,8,9,10\}$ & 4  & $7 - 4 = 3$ \\
7 & 11 & $\{11,12,13,\ldots\}$ & $\ge 3$ & $11 - 7 = 4$ \\
\bottomrule
\end{tabular}
\end{center}

The gaps between successive jumps are $3$ and $4$.  If this pattern
continues, the next jump would occur at $k = 11 + 5 = 16$.

\begin{conjecture}[Staircase growth]
\label{conj:staircase}
The chromatic number $\chi(\GZk)$ grows without bound as $k \to \infty$,
with jumps occurring at $k = 4, 7, 11, 16, 22, \ldots$, where the gaps
between successive jump points form the sequence $3, 4, 5, 6, \ldots$
(i.e.\ the $n$-th jump occurs at $k = \binom{n+2}{2} + 1$ for $n \ge 1$,
counting from the $\chi = 5$ jump).
\end{conjecture}

Quick, lighter-budget TabuCol runs through $k = 20$ suggest that $7$~colours
continue to suffice and $6$~colours continue to fail, consistent with the
next jump not occurring until at least $k = 16$, but these preliminary runs
have not been checked as rigorously as $k = 11$--$13$.


%======================================================================
\section{Chromatic Polynomial}
\label{sec:chrompoly}
%======================================================================

\subsection{Full polynomial for $k = 3$}

The chromatic polynomial of $G(\Z_3)$, computed exactly by
deletion--contraction, is:
\[
  P(G(\Z_3),\, t) \;=\; t^9 - 21t^8 + 196t^7 - 1061t^6 + 3625t^5
    - 7932t^4 + 10722t^3 - 8062t^2 + 2532t.
\]
Its leading coefficients confirm the structural invariants: the coefficient
of~$t^8$ is $-|E| = -21$, and the coefficient of~$t^7$ is
$\binom{|E|}{2} - T = \binom{21}{2} - 14 = 210 - 14 = 196$, consistent
with $T(3) = 14$ triangles (Table~\ref{tab:triangles}).  The smallest
positive integer root is $t = 4$, confirming $\chi = 4$.

Sample values: $P(4) = 48$, $P(5) = 6360$, $P(6) = 120\,240$,
$P(7) = 1\,062\,600$.


\subsection{Partial values for $k = 4$}

The full chromatic polynomial of $G(\Z_4)$ has degree~$16$ and has not been
computed (the graph is too dense for plain deletion--contraction at this
size).  Individual values found by exhaustive colouring counts are:
\[
  P(G(\Z_4),\, t) = 0 \;\text{ for } t = 0, 1, 2, 3, 4;
  \qquad P(5) = 41\,040;\qquad P(6) = 25\,536\,960.
\]
That $P(4) = 0$ and $P(5) > 0$ confirms $\chi(G(\Z_4)) = 5$.


\subsection{Whitney coefficients from structural invariants}

By Whitney's broken-circuit theorem, the first three coefficients of the
chromatic polynomial $P(G, t) = \sum_{i=0}^{n} (-1)^{n-i} a_{n-i}\, t^i$
of any graph on $n$~vertices are determined by the edge count and triangle
count:
\begin{align}
  a_0 &= 1, \nonumber \\
  a_1 &= |E|, \label{eq:a1} \\
  a_2 &= \binom{|E|}{2} - T. \label{eq:a2}
\end{align}
Substituting the formulas $|E| = k(k{-}1)(2k{+}1)/2$ (proved) and
$T = k(k{-}1)(4k{-}5)/3$ (conjectured), we obtain:
\begin{align*}
  a_1(\GZk) &= \frac{k(k-1)(2k+1)}{2}, \\[4pt]
  a_2(\GZk) &= \binom{k(k{-}1)(2k{+}1)/2}{2} - \frac{k(k-1)(4k-5)}{3}.
\end{align*}
These give the first nontrivial coefficients of $P(\GZk, t)$ in closed form
for all $k \ge 3$ (with $a_2$ conditional on Conjecture~\ref{conj:triangles}).
The next coefficient $a_3$ depends on the $4$-cycle count, which we have not
investigated.


%======================================================================
\section{Properties Proved for All $k \ge 3$}
\label{sec:summary}
%======================================================================

For reference, we collect the results that hold for every $k \ge 3$ (not
just the computed range $3 \le k \le 13$):

\begin{enumerate}
  \item \textbf{Order:} $|V| = k^2$. \hfill (Definition)
  \item \textbf{Size:} $|E| = k(k{-}1)(2k{+}1)/2$. \hfill (Theorem~\ref{thm:edges})
  \item \textbf{Degree sequence:} $k$ vertices of degree $2k{-}2$, and
    $k(k{-}1)$ vertices of degree $2k{-}1$. \hfill (Proposition~\ref{prop:deg})
  \item \textbf{Connectivity:} connected, with diameter exactly~$2$.
    \hfill (Theorem~\ref{thm:diam})
  \item \textbf{Clique number:} $\omega = 3$. \hfill (Theorem~\ref{thm:clique})
  \item \textbf{Self-loops:} exactly at $(p, 0)$ for $p \in \Zk$.
    \hfill (Remark~\ref{rem:selfloops})
  \item \textbf{Chromatic number:} $\chi \ge 4$ (since
    $\omega = 3$ and a greedy argument on the dense neighbourhood structure
    requires at least~$4$).
\end{enumerate}


%======================================================================
\section{Open Problems}
\label{sec:open}
%======================================================================

\begin{enumerate}
  \item \textbf{Chromatic number for $k \ge 14$.}  Does $\chi$ stay at~$7$
    or eventually jump to~$8$?  Preliminary TabuCol runs suggest $\chi = 7$
    persists through $k = 20$, but these have not been verified rigorously.

  \item \textbf{Rigorous lower bound for $k = 11$--$13$.}  Turn the heuristic
    evidence that $\GZk$ is not $6$-colourable for $k = 11, 12, 13$ into a
    formal proof, either via a SAT/UNSAT certificate or a structural argument.

  \item \textbf{Prove the staircase pattern.}  Explain structurally why the
    jumps in~$\chi$ occur at $k = 4, 7, 11, \ldots$ when the clique number
    is only~$3$ throughout.  In particular, is
    Conjecture~\ref{conj:staircase} correct?

  \item \textbf{Independence number.}  Prove
    Conjecture~\ref{conj:alpha}: $\alpha(\GZk) = \lfloor k^2/4 \rfloor$.

  \item \textbf{Triangle count.}  Prove
    Conjecture~\ref{conj:triangles}: $T(\GZk) = k(k{-}1)(4k{-}5)/3$.

  \item \textbf{Chromatic polynomial.}  Compute the full chromatic polynomial
    of $G(\Z_4)$ (degree~$16$) and, if feasible, $G(\Z_5)$ (degree~$25$).
    Determine whether the higher Whitney coefficients also have closed forms.
\end{enumerate}


%======================================================================
% References
%======================================================================
\begin{thebibliography}{9}

\bibitem{hertz1987}
A.~Hertz and D.~de~Werra,
\emph{Using tabu search techniques for graph coloring},
Computing \textbf{39} (1987), 345--351.

\bibitem{rawat2026}
R.~Rawat and A.~Sehgal,
\emph{Central digraphs and their chromatic properties},
2026.
% TODO: Add full citation once published.

\end{thebibliography}

\end{document}
